\documentclass[../../../include/open-logic-section]{subfiles}

\begin{document}

\olfileid{sth}{z}{arbintersections}
\olsection{પરિશિષ્ટ: સંવરણ, ગણરચના અને છેદ}

\olref[sth][z][milestone]{sec}માં આપણે
\olref[sfr][][]{part}ના સહજ રીતે કરેલા કામ તરફ પાછા જોઈને
તે~$\Zminus$માં કરી શકાય છે તે તપાસવાનું સૂચવ્યું હતું.
તે સલાહ અનુસરી હોય, તો કદાચ \emph{એક} મુદ્દે તમે અટક્યા
હશો: ડેડેકિન્ડની \emph{સંવરણો}ની ચર્ચામાં \emph{છેદ}નો ઉપયોગ.

\olref[sfr][infinite][dedekind]{Closure} પરથી યાદ કરો કે
\begin{align*}
  \closureofunder{f}{o} & = \bigcap\Setabs{X}{o \in X \text{ અને $X$ એ
$f$-સંવૃત છે}}.
\intertext{આનો સામાન્ય આકાર નીચેના સ્વરૂપની વ્યાખ્યા છે:}
  C & = \bigcap\Setabs{X}{\phi(X)}.
\end{align*}
પણ અહીં ચેતવું જોઈએ: સહજ ગણરચના નિષ્ફળ જાય છે, તેથી
$\Setabs{X}{\phi(X)}$ અસ્તિત્વ ધરાવે તેની ખાતરી નથી.
આવી વ્યાખ્યાઓ આમ જોખમી રીતે \emph{છેતરપિંડી} જેવી લાગે છે.

સદભાગ્યે તે છેતરપિંડી નથી; અથવા, જે સ્વરૂપે તે છે તેમાં
છેતરપિંડી \emph{હોય}, તો તેને યોગ્ય બનાવવા થોડું પ્રામાણિક
કામ કરી શકાય છે. આ પ્રામાણિક કામની પૂર્વઝાંખી
\olref[sth][z][sep]{prop:intersectionsexist}માં આપી હતી,
જ્યારે દરેક $A \neq \emptyset$ માટે $\bigcap A$ શા માટે
અસ્તિત્વ ધરાવે છે તે સમજાવ્યું હતું. પણ હવે તેને સ્પષ્ટ
રીતે રજૂ કરીશું.

ઘટકો દ્વારા નિર્ધારિત સમાનતા અપાયેલી હોય ત્યારે $C$ને
$\bigcap\Setabs{X}{\phi(X)}$ તરીકે વ્યાખ્યાયિત કરવાનો પ્રયાસ
કરીએ, તો હકીકતમાં આપણે નીચેની શરત પાળતી વસ્તુ $C$
જ માગીએ છીએ:
\begin{equation}\ollabel{bicondelimarbintersection}
	\forall x(x \in C \liff \forall X(\phi(X) \lif x \in X))\tag{*}
\end{equation}
હવે ધારો કે એવો \emph{કોઈ} ગણ $S$ છે કે $\phi(S)$.
ત્યારે \olref{bicondelimarbintersection} મેળવવા માટે
\emph{પૃથક્કરણ}થી $C$ને સીધો નીચે પ્રમાણે વ્યાખ્યાયિત
કરી શકીએ છીએ:
\[
	C = \Setabs{x \in S}{\forall X(\phi(X) \lif x \in X)}.
\]
આ વ્યાખ્યાથી ઇચ્છિત \olref{bicondelimarbintersection}
મળે છે તે તપાસવાનું સ્વાધ્યાય તરીકે છોડીએ છીએ.
%  fix $x$. First suppose that $x \in C$, as (re)defined; then both $x
%  \in S$ and $\forall X(\phi(X) \lif x \in X)$; but since
%  $\phi(S)$, this reduces to the condition that $\forall X(\phi(X)
%  \lif x \in X)$. Conversely, suppose $\forall X(\phi(X) \lif
%  x \in X)$; then, since $\phi(S)$, we also have that $x \in S$; so
%  $x \in c$, as (re)defined. 
આ સામાન્ય રીતથી છેદોની વ્યાખ્યામાં દેખાતા સહજ ગણરચનાના
કોઈ પણ ઉપયોગને ટાળી શકાય છે. આ વિચારણા શરૂ કરાવનાર
ખાસ કિસ્સામાં, એટલે કે $\closureofunder{f}{o}$ માટે,
આ રીત કેવી રીતે કામ કરે તે જોઈએ. આપણે
\olref[sfr][infinite][dedekind]{closureproperties}ની
સાબિતી $o \in \ran{f}\cup\{o\}$ અને
$\ran{f} \cup \{o\}$ એ $f$-સંવૃત છે તે નોંધીને શરૂ
કરી હતી. તેથી ઇચ્છિત વસ્તુ આ રીતે વ્યાખ્યાયિત કરી શકીએ:
\[
	\closureofunder{f}{o} = \Setabs{x \in \ran{f} \cup \{o\}}{(\forall X \ni o)(X \text{ એ $f$-સંવૃત છે} \lif x \in X)}.
\]

\end{document}
